\section{A complete before-and-after formalization}
\label{sec:latex-style-before-after}

Consider the sentence ``Clearly, if every \(P\) is a \(Q\) and something is a
\(P\), then something is a \(Q\).''  As written, it does not identify a
signature, a calculus, the status of the letters, or the inference that uses the
existential premise.  It is therefore not yet a formal-logical claim.

Fix a first-order signature with unary relation symbols \(P,Q\), use the
natural-deduction calculus \(\ND_{\FOL}\), and let \(x,y\) be object-language
variables.  The corrected claim is
\[
  \forall x(P(x)\to Q(x)),\ \exists xP(x)
  \derives_{\ND_{\FOL}} \exists xQ(x).
\]
Here is a complete derivation.  The superscript \(i\) marks the assumption
discharged by \(\exists E\):
\[
\begin{array}{rll}
1.&\forall x(P(x)\to Q(x))&\text{premise},\\
2.&\exists xP(x)&\text{premise},\\
3.&[P(y)]^i&\text{assumption},\\
4.&P(y)\to Q(y)&1,\ \forall E,\\
5.&Q(y)&3,4,\ \to E,\\
6.&\exists xQ(x)&5,\ \exists I,\\
7.&\exists xQ(x)&2,\ 3\text{--}6,\ \exists E\ (i).
\end{array}
\]
Choose \(y\) absent from the conclusion, from every other open assumption, and
from the existential premise.  Those are exactly the active
\(\exists E\) side conditions.  Lines 4--6 show that an arbitrary witness for
the premise supplies a witness for the conclusion; line 7 discharges the only
open assumption.  The revised passage therefore fixes the language,
calculus, metavariable standing, rule applications, discharge, and freshness
condition that the original sentence omitted.

